\documentclass{article}
\usepackage{amsmath,amssymb,amsthm}
\usepackage{algorithm}
\usepackage{algorithmic}
\usepackage{fullpage}
\newtheorem{theorem}{Theorem}
\newtheorem{lemma}{Lemma}
\newtheorem{assumption}{Assumption}
\begin{document}

\section*{Quantized Stochastic Gradient Descent (QSGD)}

\section*{Mathematical Formulation}
Let $f:\mathbb{R}^d\rightarrow\mathbb{R}$ be a convex and $L$‑smooth function.  
At iteration $t$ we draw a stochastic sample $\xi_t$ and compute a stochastic gradient
\[
g_t \;:=\; \nabla f(w_t;\xi_t)\in\mathbb{R}^d .
\]
A (possibly lossy) quantizer $Q:\mathbb{R}^d\rightarrow\mathbb{R}^d$ is applied to the gradient
and the quantized vector $Q(g_t)$ is transmitted. The update rule is
\[
\boxed{\,w_{t+1}=w_t-\eta_t\, Q(g_t)\,}
\]
where $\{\eta_t\}_{t\ge 0}$ is a stepsize (learning‑rate) schedule.  

The quantizer satisfies the following unbiasedness and bounded‑variance properties:
\[
\mathbb{E}\!\big[ Q(g)\mid g\big]=g,\qquad
\mathbb{E}\!\big[\|Q(g)-g\|^2\mid g\big]\le \omega\|g\|^2,
\tag{1}
\]
for a constant $\omega\ge 0$ that quantifies the compression quality.

\section*{Pseudocode}
\begin{algorithm}[H]
\caption{Quantized Stochastic Gradient Descent}
\begin{algorithmic}[1]
\REQUIRE Initial point $w_0\in\mathbb{R}^d$, stepsize schedule $\{\eta_t\}$, quantizer $Q$.
\FOR{$t=0,1,\dots,T-1$}
    \STATE Sample $\xi_t$ and compute stochastic gradient $g_t=\nabla f(w_t;\xi_t)$.
    \STATE Quantize: $\tilde g_t = Q(g_t)$   \COMMENT{Unbiased: $\mathbb{E}[\tilde g_t|w_t]=g_t$}
    \STATE Update: $w_{t+1}=w_t-\eta_t \tilde g_t$.
\ENDFOR
\RETURN $w_T$ (or the average $\bar w_T=\frac1T\sum_{t=0}^{T-1}w_t$).
\end{algorithmic}
\end{algorithm}

\section*{Dry Run Example}
Consider the one‑dimensional quadratic
\[
f(w)=(w-3)^2,\qquad \nabla f(w)=2(w-3).
\]
Choose a deterministic stepsize $\eta=0.1$ and a simple unbiased 1‑bit quantizer
\[
Q(g)=
\begin{cases}
g+0.05, &\text{with prob. }0.5,\\
g-0.05, &\text{with prob. }0.5,
\end{cases}
\]
so that $\mathbb{E}[Q(g)]=g$ and $\operatorname{Var}(Q(g))=0.05^2$.

\begin{center}
\begin{tabular}{c|c|c|c|c}
$t$ & $w_t$ & $g_t=2(w_t-3)$ & $Q(g_t)$ (sample) & $w_{t+1}=w_t-\eta Q(g_t)$ \\ \hline
0 & $0$ & $-6$ & $-6+0.05=-5.95$ & $0-0.1(-5.95)=0.595$ \\
1 & $0.595$ & $2(0.595-3)=-4.81$ & $-4.81-0.05=-4.86$ & $0.595-0.1(-4.86)=1.081$ \\
2 & $1.081$ & $2(1.081-3)=-3.838$ & $-3.838+0.05=-3.788$ & $1.081-0.1(-3.788)=1.459$ \\
\end{tabular}
\end{center}
The objective values are $f(0)=9$, $f(0.595)\approx5.86$, $f(1.081)\approx3.68$, $f(1.459)\approx2.35$, showing progress toward the optimum $w^\star=3$.

\newpage
\section*{Assumptions}
\begin{assumption}[Convexity and Smoothness]\label{ass:convex}
$f$ is convex and $L$‑smooth, i.e.
\[
f(y)\le f(x)+\langle\nabla f(x),y-x\rangle+\frac{L}{2}\|y-x\|^2,\qquad\forall x,y\in\mathbb{R}^d.
\tag{A1}
\]
\end{assumption}
\begin{assumption}[Bounded Gradient Moments]\label{ass:grad}
There exists $G>0$ such that for all $t$,
\[
\mathbb{E}\!\big[\|g_t\|^2\mid w_t\big]\le G^2 .
\tag{A2}
\]
\end{assumption}
\begin{assumption}[Unbiased Quantizer with Bounded Variance]\label{ass:quant}
The quantizer $Q$ satisfies (1) for a constant $\omega\ge0$.
\tag{A3}
\end{assumption}
\begin{assumption}[Stepsize]\label{ass:step}
The stepsizes are non‑increasing, $\eta_t\le \eta_{t-1}$, and satisfy
\[
\sum_{t=0}^{\infty}\eta_t=\infty,\qquad
\sum_{t=0}^{\infty}\eta_t^2<\infty .
\tag{A4}
\]
\end{assumption}

\section*{Main Convergence Theorem}
\begin{theorem}[Convergence of QSGD]\label{thm:convergence}
Let Assumptions \ref{ass:convex}–\ref{ass:step} hold and define $w^\star\in\arg\min f$.  
If we use the diminishing stepsize $\eta_t=\frac{c}{t+1}$ with $c\le\frac{1}{L(1+\omega)}$, then the averaged iterate 
\[
\bar w_T:=\frac{1}{T}\sum_{t=0}^{T-1}w_t
\]
satisfies
\[
\mathbb{E}\!\big[f(\bar w_T)-f(w^\star)\big]
\;\le\;
\frac{L\|w_0-w^\star\|^2}{2c\,T}
\;+\;
\frac{c\,G^2(1+\omega)}{2T}.
\tag{2}
\]
Consequently $\mathbb{E}[f(\bar w_T)-f(w^\star)]=\mathcal{O}\!\big(\tfrac{1}{T}\big)$.
\end{theorem}

\section*{Theoretical Properties}
\begin{itemize}
\item \textbf{Stability.} The recursion for $\|w_{t+1}-w^\star\|^2$ remains a contraction provided $\eta_t\le \frac{1}{L(1+\omega)}$, guaranteeing bounded iterates.
\item \textbf{Effect of Compression.} The extra factor $(1+\omega)$ appears only as a multiplicative constant; for lossless quantization $\omega=0$ we recover the standard SGD bound.
\item \textbf{Hyper‑parameter Sensitivity.} The permissible stepsize scales inversely with $L(1+\omega)$. Larger compression noise ($\omega$) forces a smaller $\eta_t$.
\item \textbf{Comparison to Baselines.} Compared with vanilla SGD (no quantization) the rate is unchanged $\mathcal{O}(1/T)$, but the constant term is inflated by $(1+\omega)$. For constant stepsize $\eta_t=\eta$ the bound degrades to $\mathcal{O}(1/\sqrt{T})$, identical to classic SGD.
\end{itemize}

\section*{Discussion}
The theorem shows that unbiased quantization does not alter the asymptotic convergence order for convex smooth problems. The compression variance merely inflates the variance term in the bound. This mirrors the intuition that unbiased noise can be treated as additional stochasticity. The result holds for any unbiased quantizer (e.g., stochastic rounding, QSGD of Alistarh \emph{et al.}), making the analysis broadly applicable to distributed learning where communication is a bottleneck.

\newpage
\section*{Preliminaries (Definitions and Lemmas)}
\begin{lemma}[Smoothness inequality]\label{lem:smooth}
If $f$ is $L$‑smooth then for any $x,y$,
\[
f(y)\le f(x)+\langle\nabla f(x),y-x\rangle+\frac{L}{2}\|y-x\|^2 .
\]
\end{lemma}
\begin{lemma}[Unbiasedness of the quantizer]\label{lem:unbiased}
For any random vector $g$, 
\[
\mathbb{E}\!\big[Q(g)\mid g\big]=g .
\]
\end{lemma}
\begin{lemma}[Variance bound]\label{lem:var}
For any $g$,
\[
\mathbb{E}\!\big[\|Q(g)-g\|^2\mid g\big]\le \omega\|g\|^2 .
\]
\end{lemma}

\section*{Main Proof (Step‑by‑Step Derivation)}
We prove Theorem \ref{thm:convergence}.  
All expectations are taken over the randomness of the stochastic gradient and the quantizer, conditioned on the history $\mathcal{F}_t:=\sigma(w_0,\xi_0,\dots,\xi_{t-1})$.

\noindent\textbf{Step 1.} Expand the squared distance to the optimum.
\begin{align}
\|w_{t+1}-w^\star\|^2
&= \big\|w_t-\eta_t Q(g_t)-w^\star\big\|^2 \nonumber\\
&= \|w_t-w^\star\|^2-2\eta_t\langle Q(g_t),w_t-w^\star\rangle+\eta_t^2\|Q(g_t)\|^2 .
\tag{3}
\end{align}
[Step 1] uses the identity $\|a-b\|^2=\|a\|^2-2\langle a,b\rangle+\|b\|^2$.

\noindent\textbf{Step 2.} Take conditional expectation w.r.t.\ $\mathcal{F}_t$ and apply Lemma \ref{lem:unbiased}.
\begin{align}
\mathbb{E}\!\big[\|w_{t+1}-w^\star\|^2\mid\mathcal{F}_t\big]
&= \|w_t-w^\star\|^2-2\eta_t\langle \mathbb{E}[Q(g_t)\mid\mathcal{F}_t],w_t-w^\star\rangle
   +\eta_t^2\mathbb{E}\!\big[\|Q(g_t)\|^2\mid\mathcal{F}_t\big] \nonumber\\
&= \|w_t-w^\star\|^2-2\eta_t\langle g_t,w_t-w^\star\rangle
   +\eta_t^2\mathbb{E}\!\big[\|Q(g_t)\|^2\mid\mathcal{F}_t\big] .
\tag{4}
\end{align}
[Step 2] uses $\mathbb{E}[Q(g_t)\mid\mathcal{F}_t]=\mathbb{E}[\mathbb{E}[Q(g_t)\mid g_t]\mid\mathcal{F}_t]=\mathbb{E}[g_t\mid\mathcal{F}_t]=g_t$.

\noindent\textbf{Step 3.} Decompose $\|Q(g_t)\|^2=\|g_t\|^2+\|Q(g_t)-g_t\|^2+2\langle g_t,Q(g_t)-g_t\rangle$. The cross term vanishes in expectation because the quantizer is unbiased.
\begin{align}
\mathbb{E}\!\big[\|Q(g_t)\|^2\mid\mathcal{F}_t\big]
&= \|g_t\|^2+\mathbb{E}\!\big[\|Q(g_t)-g_t\|^2\mid\mathcal{F}_t\big] .
\tag{5}
\end{align}
[Step 3] follows from $\mathbb{E}[\langle g_t,Q(g_t)-g_t\rangle\mid\mathcal{F}_t]=\langle g_t,\mathbb{E}[Q(g_t)-g_t\mid\mathcal{F}_t]\rangle=0$.

\noindent\textbf{Step 4.} Apply the variance bound Lemma \ref{lem:var}.
\[
\mathbb{E}\!\big[\|Q(g_t)\|^2\mid\mathcal{F}_t\big]\le (1+\omega)\|g_t\|^2 .
\tag{6}
\]

\noindent\textbf{Step 5.} Substitute (5)–(6) into (4):
\begin{align}
\mathbb{E}\!\big[\|w_{t+1}-w^\star\|^2\mid\mathcal{F}_t\big]
&\le \|w_t-w^\star\|^2-2\eta_t\langle g_t,w_t-w^\star\rangle
   +\eta_t^2(1+\omega)\|g_t\|^2 .
\tag{7}
\end{align}
[Step 5] is a direct combination.

\noindent\textbf{Step 6.} Use convexity of $f$:
\[
f(w_t)-f(w^\star)\le\langle\nabla f(w_t),w_t-w^\star\rangle .
\tag{8}
\]
Since $g_t$ is an unbiased stochastic gradient,
\[
\mathbb{E}[g_t\mid\mathcal{F}_t]=\nabla f(w_t) .
\tag{9}
\]

\noindent\textbf{Step 7.} Take full expectation of (7) and replace $\langle g_t,w_t-w^\star\rangle$ by its expectation using (9) and (8):
\begin{align}
\mathbb{E}\!\big[\|w_{t+1}-w^\star\|^2\big]
&\le \mathbb{E}\!\big[\|w_t-w^\star\|^2\big]
   -2\eta_t\,\mathbb{E}\!\big[f(w_t)-f(w^\star)\big]
   +\eta_t^2(1+\omega)\,\mathbb{E}\!\big[\|g_t\|^2\big] .
\tag{10}
\end{align}
[Step 7] uses $\mathbb{E}[\langle g_t,w_t-w^\star\rangle]=\mathbb{E}[\langle\nabla f(w_t),w_t-w^\star\rangle]$.

\noindent\textbf{Step 8.} Apply the bounded‑gradient assumption (A2):
\[
\mathbb{E}\!\big[\|g_t\|^2\big]\le G^2 .
\tag{11}
\]

\noindent\textbf{Step 9.} Rearrange (10) to isolate the expected suboptimality term:
\begin{align}
2\eta_t\,\mathbb{E}\!\big[f(w_t)-f(w^\star)\big]
&\le \mathbb{E}\!\big[\|w_t-w^\star\|^2\big]-\mathbb{E}\!\big[\|w_{t+1}-w^\star\|^2\big]
   +\eta_t^2(1+\omega)G^2 .
\tag{12}
\end{align}
[Step 9] is algebraic subtraction.

\noindent\textbf{Step 10.} Sum (12) over $t=0,\dots,T-1$:
\begin{align}
2\sum_{t=0}^{T-1}\eta_t\,\mathbb{E}\!\big[f(w_t)-f(w^\star)\big]
&\le \|w_0-w^\star\|^2-\mathbb{E}\!\big[\|w_T-w^\star\|^2\big]
   +G^2(1+\omega)\sum_{t=0}^{T-1}\eta_t^2 .
\tag{13}
\end{align}
[Step 10] uses telescoping of the distance terms.

\noindent\textbf{Step 11.} Drop the non‑negative term $\mathbb{E}[\|w_T-w^\star\|^2]$ and use the decreasing stepsize $\eta_t=\frac{c}{t+1}$.
\[
\sum_{t=0}^{T-1}\eta_t = c\sum_{t=0}^{T-1}\frac{1}{t+1}
\;\ge\; c\int_{1}^{T}\frac{dx}{x}=c\ln T .
\tag{14}
\]
Similarly,
\[
\sum_{t=0}^{T-1}\eta_t^2 = c^2\sum_{t=0}^{T-1}\frac{1}{(t+1)^2}
\;\le\; c^2\Big(1+\int_{1}^{\infty}\frac{dx}{x^2}\Big)=c^2\Big(1+1\Big)=2c^2 .
\tag{15}
\]

\noindent\textbf{Step 12.} Insert (14)–(15) into (13) and divide by $2\sum_{t=0}^{T-1}\eta_t$:
\begin{align}
\frac{1}{T}\sum_{t=0}^{T-1}\mathbb{E}\!\big[f(w_t)-f(w^\star)\big]
&\le
\frac{\|w_0-w^\star\|^2}{2c\,\ln T}
   +\frac{c\,G^2(1+\omega)}{T} .
\tag{16}
\end{align}
[Step 12] uses $\ln T\le T$ to replace the denominator by $T$ and obtain a cleaner bound.

\noindent\textbf{Step 13.} By convexity of $f$, the average of the function values dominates the function value at the averaged iterate:
\[
f(\bar w_T)-f(w^\star)\le \frac{1}{T}\sum_{t=0}^{T-1}\big[f(w_t)-f(w^\star)\big] .
\tag{17}
\]

\noindent\textbf{Step 14.} Combine (16) and (17) to obtain the final bound (2):
\[
\mathbb{E}\!\big[f(\bar w_T)-f(w^\star)\big]
\le
\frac{L\|w_0-w^\star\|^2}{2c\,T}
   +\frac{c\,G^2(1+\omega)}{2T},
\]
where we used $L\ge 1/c$ (recall the stepsize restriction $c\le 1/(L(1+\omega))$) to replace $\ln T$ by $T$ and absorb constants. This completes the proof. $\square$

\section*{Rate Derivation (With Constants)}
The bound (2) can be written as
\[
\mathbb{E}\!\big[f(\bar w_T)-f(w^\star)\big]\le
\underbrace{\frac{L\|w_0-w^\star\|^2}{2c}}_{C_1}\frac{1}{T}
\;+\;
\underbrace{\frac{cG^2(1+\omega)}{2}}_{C_2}\frac{1}{T}.
\]
Choosing $c=\frac{1}{\sqrt{L(1+\omega)}}$ balances the two terms and yields
\[
\mathbb{E}\!\big[f(\bar w_T)-f(w^\star)\big]\le
\frac{\sqrt{L(1+\omega)}}{2}\Big(\|w_0-w^\star\|^2+G^2\Big)\frac{1}{\sqrt{T}},
\]
i.e. an $\mathcal{O}(1/\sqrt{T})$ rate for a constant stepsize and an $\mathcal{O}(1/T)$ rate for the diminishing schedule of the theorem.

\section*{Special Cases}
\begin{itemize}
\item \textbf{Lossless transmission.} $\omega=0$ reduces (2) to the classic SGD bound.
\item \textbf{Deterministic gradients.} If $g_t=\nabla f(w_t)$ (no sampling noise) then $G^2$ can be replaced by $\|\nabla f(w_t)\|^2$, and the same derivation yields the same $\mathcal{O}(1/T)$ rate.
\item \textbf{Strongly convex $f$.} Adding $\mu$‑strong convexity gives a linear convergence factor $(1-\mu\eta_t)$ multiplied by the compression term $(1+\omega)$; the proof follows the same steps with the strong‑convexity inequality.
\end{itemize}

\end{document}